% Français 9115, batch 8 — Gallica views 141–160.
% Pass: Opus 5 (claude-opus-5), 2026-09-15 — first pass, unchecked against
% the views by a human.
%
% What the twenty views hold. Two things, unequal in length, and neither
% of them Germain's own argument. Views 140 to 143 are small leaves mounted
% on larger support sheets: view 140 (batch 7) is a note on the element
% $dx\,dy\sqrt{1+p^2+q^2}$ of a surface, and view 141 is its blank verso.
% View 142 is another small leaf, headed « Gauss courbure des surfaces »:
% lines in French, every line opened with inverted commas, rendering
% the opening of a text on the representation of directions by points of a
% sphere of radius 1 — the direction of a line by the point where the
% parallel radius ends, the three axes by three points a quadrant apart, the
% position of a plane by the great circle of the parallel plane. The leaf's
% own title names Gauss; nothing further is said here about which text it
% renders. View 143 is its blank verso.
%
% From view 144 to the end of the batch, and without a break, comes a copy
% in the same hand as the French leaves of a Latin memoir whose first leaf
% gives its own source: « Acta academiæ scientiarum imperialis
% Petropolitanæ pro anno 1779 Pars prior », « Investigatio motuum quibus
% laminæ et virgæ elasticæ contremiscunt auctore L. Eulero ». It is a fair
% copy, written on both sides of each leaf, and each view carries at its
% head a number in parentheses, (103) at view 144 running to (119) at view
% 160, one to a view without a gap — very probably the pages of the printed
% volume being copied, which is an inference of this pass. Seventeen views,
% seventeen numbers.
%
%   v. 144–146  The preamble, why the subject is taken up again; Lemma 1, the
%               equilibrium of an elastic rod $EYF$ under any forces
%               (« Tab. II fig. 3 »), with the moment of elasticity
%               $V(dy\,ddx - dx\,ddy)/ds^3$ and the tension; Scholion § 2 on a
%               rod naturally curved, radius of osculation $\varsigma$;
%               Lemma 2, the motion of the same rod.
%   v. 147      Lemma 2 solved by d'Alembert's principle: $P'$ and $Q'$, the
%               equation of motion and the tension. Problema 1 begins
%               (« Tab. II fig. 4 »).
%   v. 148–150  Problema 1: a straight rod of uniform thickness and
%               elasticity, vibrating in any way. Masses and forces measured
%               by volumes, $V = bc^4$; differentiating twice gives
%               $E(ddy/ds^2) + \frac{cc}{2g}(ddy/dt^2) = -bc^4(d^4y/ds^4)$.
%               Corollary I, § 5: with no force, the fourth-order equation
%               whose integral « nullo adhuc modo inveniri potuisse ».
%               Corollary II, § 6: a compressing force, and the columns.
%   v. 151      Corollary II continued: a stretching weight $P$ hung over a
%               pulley (« tab. II fig. 5 »), the rod's sound a mean between an
%               elastic rod's and a string's; Corollary III, § 7, the two
%               limits $b = 0$ and $h = 0$.
%   v. 152      Scholion § 8: longitudinal forces set aside; the three
%               equations I, II, III kept for the rods fixed by pins.
%   v. 153–156  Problema § 9: the regular vibrations, isochronous with a
%               simple pendulum of length $k$. $y = e^{\lambda s}$,
%               $\lambda^4 = 1/f^4$ with $bcck = f^4$; the complete integral
%               with $\alpha, \beta, \gamma, \delta$, sines and cosines; time
%               of a vibration and number per second. Corollary 1, § 10, the
%               series; Corollary 2, § 11, the derivatives in $s$; Corollary
%               3, § 12, velocity and acceleration; Scholion I, § 13, begins.
%   v. 157–158  Scholion I continued: three states of an end — free, simply
%               fixed, « infixus » — each giving two conditions. Scholion 2,
%               § 14: six cases.
%   v. 159–160  The table of the six cases; « Evolutio casus I quo virgae
%               elasticae uterque terminus prorsus est liber », Problema
%               § 15, and its solution as far as the equation
%               $\cos\omega = 2/(e^{\omega} + e^{-\omega})$ and the root
%               $\omega = 0$, a rod at rest. The copy runs on into batch 9.
%
% The numbers on the leaves, and why no \folio{} is written. The ink series
% of the earlier batches runs on at the top right of the leaves that face
% up: 70 at view 142, then 71, 72, 73, 74, 75, 76, 77, 78, 79 at views 144,
% 146, 148, 150, 152, 154, 156, 158, 160 (69 stands at view 140, in batch
% 7). The odd views from 145 to 159 carry no ink number, only the
% parenthesised page number of the copy. The copy's numbers (103)–(119) are
% a third kind, recorded in the notes at the head of each view. View 142
% also carries « 1 » and « 2 » in its left margin, numbering its two
% paragraphs. On the mounting sheets, at the top right of the rectos (and
% showing through at the top left of the versos), three short oblique
% pencil strokes stand at the same place on every sheet; they are not a
% figure and are not read as foliation. No pencilled foliation of the
% Bibliothèque was read anywhere in the batch, and no \folio{} is written,
% as in batches 1 to 5. Every number above was read on its view; none is
% inferred.
%
% Notation. The copy's letters are kept: $V$ the absolute elasticity, $S$
% the mass per unit length, $P$, $Q$ and $P'$, $Q'$ the forces, $g$ the
% fall in one second, parentheses « uncinulis » for partial differentials
% (the copy's own word, view 147), « sin. » and « cos. » with their stops
% set as \text{sin.} and \cos. where written. The radius of osculation of
% the naturally curved rod is a looped Greek letter, set $\varsigma$; the
% arbitrary coefficient of the general solution is a calligraphic capital,
% set $\mathfrak{M}$; the exponent letter, a small crossed letter, is set
% $\lambda$. The copy's $\alpha$ looks like « u » and its $\gamma$ like
% « v »; they are read as Greek throughout. « \&c. » renders a looped sign
% close to $\infty$ at the end of series. The ligature « æ », written in
% some words and not in others (« laminæ » but « virgae »), is kept where
% the leaf writes it. The copyist's underlining of single letters (« s »,
% « t », « f », « k ») is set \emph{}. Signs in the calculations are as
% read, and where two readings of a letter conflict with the calculation
% (view 149, the three « alteram » and $EE$, $EF$, $Ef$; view 150, $bc^2$
% beside $bc^4$; view 160, $p - q = -2\cos\omega$) the leaf's reading
% stands and a note says so.
%
% The hand. The copy is a regular, upright, fast fair hand, the hand of
% the French leaves; its Latin is easy to read at band scale. The Greek
% letters and the partial-differential denominators are the hard part, and
% were read from tight crops. The Gauss leaf is the same hand, looser.
%
% What could not be read. Nothing is marked \ill{} in this batch. Readings
% offered with doubt are marked \uncertain{}: « Gauss » and « soit » and the
% marginal « 2 » at view 142; « Illustriss. » and « que » at view 144; « a »
% at view 145; « at » at view 148; « prouti » at view 150; « clausulis » at
% view 153; « portabit » at view 157; « evolvi » and « extensione » at view
% 158.
%
% Nothing here was checked against any published transcription or edition
% of Euler's memoir or of Gauss.
\documentclass[11pt,a4paper]{article}
% The preamble shared by every transcription of the mathematicians' manuscripts in Gallica.
%
% It rests on one idea: **uncertainty compounds, it does not resolve.** A
% transcription that smooths over a doubtful word has destroyed the only thing
% it was made for — a reader must be able to tell, at every point, what was on
% the page and what was supplied. The apparatus macros below are therefore the
% critical apparatus, not decoration.
%
% The same commands are recognised by `scripts/render.mjs`, which produces the
% site's reading view, and by `scripts/tei.mjs`. A macro added here must be
% added there too, or rendering fails loudly — which is the wanted behaviour.
%
% Comments here are English, like the rest of the repository. The documents
% this preamble sets are French, and that is deliberate: see the skills.

\ProvidesPackage{germain}[2026/09/15 Transcriptions of mathematicians' manuscripts in Gallica]

\usepackage{amsmath,amssymb,amsthm}
\usepackage{mathrsfs}
% Kept from the parent preamble so the renderer's diagram support stays
% available; these manuscripts draw no commutative diagrams, and nothing here uses it.
\usepackage{tikz-cd}
\usepackage[english,french]{babel}
\usepackage{fontspec}

% --- Greek ------------------------------------------------------------------
% The text face stays fontspec's default, Latin Modern, which has no Greek:
% XeTeX drops every glyph it lacks with a line in the log and nothing on the
% page, and the Greek words the manuscripts quote vanished from the PDFs. So
% the Greek and Greek Extended blocks get a XeTeX character class of their own,
% and entering or leaving a run of it switches the family to CMU Serif — the
% same Computer Modern design, with polytonic Greek — and back. Only the
% family changes: a Greek word in \textit or \textbf keeps its shape and series.
% The transitions to and from every other class carry the tokens that class
% has towards an ordinary letter, so French spacing before « ; » or inside
% « » still holds after a Greek word. No grouping, so no brace or \endgroup
% can come out unbalanced; maths is left alone, since XeTeX inserts nothing
% there. Kept identical in `exercices.sty`.
\makeatletter
\newfontfamily\germainGreekfont{cmun}[
  Extension=.otf, NFSSFamily=germaingreek,
  UprightFont=*rm, ItalicFont=*ti, SlantedFont=*sl,
  BoldFont=*bx, BoldItalicFont=*bi, BoldSlantedFont=*bl]
\newXeTeXintercharclass\germain@greek
\def\germain@greekrange#1#2{%
  \@tempcnta=#1\relax
  \loop
    \XeTeXcharclass\@tempcnta=\germain@greek
    \ifnum\@tempcnta<#2\advance\@tempcnta\@ne\repeat}
\germain@greekrange{"0370}{"03FF}
\germain@greekrange{"1F00}{"1FFF}
\def\germain@greekprev{\rmdefault}
\def\germain@greekin{%
  \edef\germain@greekprev{\f@family}\fontfamily{germaingreek}\selectfont}
\def\germain@greekout{\fontfamily{\germain@greekprev}\selectfont}
\def\germain@greekjoin{%
  \XeTeXinterchartokenstate=\@ne
  \@tempcnta=\z@
  \loop
    \ifnum\@tempcnta=\germain@greek\else
      \XeTeXinterchartoks\germain@greek\@tempcnta=\expandafter{%
        \expandafter\germain@greekout\the\XeTeXinterchartoks\z@\@tempcnta}%
      \XeTeXinterchartoks\@tempcnta\germain@greek=\expandafter{%
        \the\XeTeXinterchartoks\@tempcnta\z@\germain@greekin}%
    \fi
    \ifnum\@tempcnta<4095\advance\@tempcnta\@ne\repeat}
% babel-french sets its own classes, and saves and restores the interchar
% state, each time a language is selected: join again after it does.
\addto\extrasfrench{\germain@greekjoin}
\addto\extrasenglish{\germain@greekjoin}
\AtBeginDocument{\germain@greekjoin}
\makeatother

\usepackage{geometry}
\geometry{a4paper,margin=25mm,marginparwidth=28mm}
\usepackage{marginnote}
\usepackage{xcolor}
\usepackage[normalem]{ulem}
\setlength{\emergencystretch}{3em}
\usepackage{hyperref}
\hypersetup{colorlinks=true,linkcolor=black,urlcolor=black,pdfborder={0 0 0}}

% --- Batch metadata ---------------------------------------------------------
% Taken as they stand by the rendering script, which shows them at the head of
% the reading view. They fix what the file claims to cover. The macro names are
% the parent project's, so that the scripts stay shared; \folder here means the
% volume's slug — `fr-9115` — and \pages counts Gallica views.

\newcommand{\folder}[1]{\gdef\grothFolder{#1}}
\newcommand{\batch}[1]{\gdef\grothBatch{#1}}
\newcommand{\pages}[2]{\gdef\grothFirst{#1}\gdef\grothLast{#2}}
\newcommand{\foldertitle}[1]{\gdef\grothFoldertitle{#1}}

% The holder's own shelfmark, printed as they write it: « Français 9115 ».
\newcommand{\shelfmark}[1]{\gdef\grothShelfmark{#1}}

% The dating, copied verbatim from the holder's catalogue — never inferred
% here. « XIXe siècle » is the BnF's; a pass that dates a leaf from its content
% says so in a \note{}, as its own inference.
\newcommand{\dating}[1]{\gdef\grothDating{#1}}

% The status notice, printed in the title block and repeated as a diagonal
% watermark on every page of the PDF. The manuscripts are in the public
% domain; what this line declares is not a right but a quality — an
% unverified first pass by a machine. The skills write the line; a file
% without it gets no watermark, so leaving it out is visible in the output.
\newcommand{\watermark}[1]{\gdef\grothWatermark{#1}}

\gdef\grothFolder{?}\gdef\grothBatch{}\gdef\grothFirst{?}\gdef\grothLast{?}%
\gdef\grothFoldertitle{}\gdef\grothShelfmark{}\gdef\grothDating{}\gdef\grothWatermark{}

% --- The résumé -------------------------------------------------------------
\newenvironment{resume}
  {\small\setlength{\parskip}{0.45em}}
  {\par\medskip\hrule\medskip}

% --- Keywords ---------------------------------------------------------------
% In English, closing the modernised reading's résumé: the modern vocabulary
% under which to search for what the leaves construct. The manifest extracts
% them to tag the volume — this line is the single source of the tags.
\newcommand{\keywords}[1]{\par\smallskip\noindent
  {\footnotesize\textsc{Keywords} --- \textit{#1}}\par}

% --- The critical apparatus -------------------------------------------------

% The Gallica view number, in the margin. This is the anchor that lets a
% disagreement over a formula be settled by looking at *one* image rather than
% a volume of 750. The site uses it to turn the facsimile as the transcription
% is scrolled. A view is one scanned image: usually one side of a leaf.
\newcommand{\page}[1]{%
  \marginnote{\footnotesize\color{gray!70}\textsf{v.\,#1}}%
  \label{page:#1}\ignorespaces}

% The BnF's pencilled foliation, where the leaf carries it and a pass can read
% it: « 198r ». This is what the literature cites — « ff. 198r–208v » — and
% the one line that turns a citation into a view. Recorded, never inferred:
% a folio a pass cannot read is not written.
\newcommand{\folio}[1]{%
  \marginnote{\footnotesize\color{gray!50}\textsf{f.\,#1}}\ignorespaces}

% The modernised reading's coarser counterpart to \page{}: which views a
% section is drawn from.
\newcommand{\pagerange}[2]{%
  \marginnote{\footnotesize\color{gray!70}\textsf{v.\,#1--#2}}%
  \ignorespaces}

% Illegible: never guessed.
\newcommand{\ill}{\textcolor{red!60!black}{[\,\ldots\,]}}

% A doubtful reading: the word is given, and flagged as uncertain.
\newcommand{\uncertain}[1]{\uline{#1}}

% An editorial addition — a missing letter, an implied word.
\newcommand{\add}[1]{\textcolor{blue!50!black}{[#1]}}

% Struck out by the writer. What was crossed out is part of the document.
\newcommand{\struck}[1]{\sout{#1}}

% The transcriber's note, on the state of the page or the mathematics.
\newcommand{\note}[1]{\par\smallskip{\small\color{gray!60!black}\textit{[#1]}}\par\smallskip}

% A marginal note of hers, distinct from the body of the text.
\newcommand{\marginal}[1]{\par\smallskip\noindent\hspace{1em}%
  {\small\color{gray!60!black}#1}\par\smallskip}

% --- Title ------------------------------------------------------------------
% The title block is French, like the documents themselves.

\AddToHook{shipout/background}{%
  \ifx\grothWatermark\empty\else
    \begin{tikzpicture}[remember picture, overlay]
      \node[rotate=52, scale=3.4, align=center, text=gray!18,
            font=\sffamily\bfseries]
        at (current page.center) {\grothWatermark};
    \end{tikzpicture}%
  \fi}

\AtBeginDocument{%
  \begin{center}
    {\large\bfseries Manuscrits de mathématiciens (Gallica) ---
      \ifx\grothShelfmark\empty\grothFolder\else\grothShelfmark\fi}\\[2pt]
    {\itshape\grothFoldertitle}\\[4pt]
    {\small vues \grothFirst--\grothLast
      \ifx\grothBatch\empty\else\ (lot \grothBatch)\fi}\\[2pt]
    \ifx\grothDating\empty\else
      {\small\color{gray!60!black}Datation du catalogue : \grothDating}\\[2pt]
    \fi
    \ifx\grothWatermark\empty\else
      {\small\bfseries\color{red!45!gray}\grothWatermark}\\[2pt]
    \fi
    {\footnotesize\color{gray!60!black}Transcription établie d'après les images IIIF de Gallica.
     Source gallica.bnf.fr / Bibliothèque nationale de France.\\
     Les lectures incertaines sont soulignées, les illisibles notées [\,\ldots\,],
     les ajouts entre crochets ; \textsf{v.} numérote les vues de Gallica,
     \textsf{f.} la foliotation au crayon de la BnF.}
  \end{center}
  \vspace{1em}}

\endinput


\folder{fr-9115}
\shelfmark{Français 9115}
\batch{8}
\pages{141}{160}
\dating{1801-1900}
\watermark{Lecture automatique\\non vérifiée}
\foldertitle{Papiers de Sophie GERMAIN. — Recueil de dissertations et problèmes mathématiques et physiques.}

\begin{document}

\note{Numérotation des feuillets. La série à l'encre des lots précédents se
poursuit en haut à droite des feuillets : 70 (vue 142), 71, 72, 73, 74, 75,
76, 77, 78, 79 aux vues 144, 146, 148, 150, 152, 154, 156, 158, 160. Les
vues impaires de 145 à 159 ne portent pas de nombre à l'encre. La copie
latine porte en outre, en tête de chaque vue de 144 à 160, un nombre entre
parenthèses, de (103) à (119), qui semble reporter la page de l'imprimé
copié : c'est une inférence. Tous ces nombres ont été lus sur la vue ;
aucun n'est déduit. Les feuilles de support portent, au même endroit sur
chacune, trois petits traits obliques au crayon qui ne forment pas un
chiffre : aucune foliotation de la Bibliothèque n'a été lue sur ces vues, et
aucun « f. » n'est donc porté ici. Les vues 141 et 143, versos blancs des
petits feuillets des vues 140 et 142, ne sont pas transcrites.}

\page{142}
\uncertain{Gauss} \emph{courbure des surfaces}

1 . . . « \struck{tous les points} la situation dans l'espace est déterminée
par trois coordonnées, \uncertain{soit} par les distances à trois plans fixes
perpendiculaires entre eux. On décrit un cercle de rayon 1 autour d'un
centre arbitraire, la direction d'une ligne sera représentée par le point de
la sphère auquel aboutit le rayon qui est parallèle à cette ligne. Si les
trois ordonnées ont leur origine au centre de la sphère elles seront les
rayons séparés par des quarts de cercle ou quadrans leur direction sera
représentée par (1) (x) (2) y (3) z.

\uncertain{2}. « la situation d'un plan peut être représentée par le plus
grand cercle formé par l'intersection de la sphère et d'un plan qui lui
seroit parallèle
\note{petit feuillet monté sur une feuille de support, écrit à l'encre ; le nombre 70 est en haut
à droite. Le titre est souligné d'un trait ; le premier mot, lu « Gauss »,
est d'une écriture plus lâche. Chaque ligne du texte, depuis « tous les
points », commence par des guillemets : c'est une citation, que les
guillemets initiaux signalent seuls ; ils ne sont pas répétés ici. Les mots
« tous les points » sont barrés, et un guillemet est rouvert après eux. Au
deuxième alinéa, « qui » est écrit au-dessus de « soit », sans que ce mot
soit clairement biffé ; « cercle », à la huitième ligne, est écrit sur un
autre mot. Les chiffres « 1 » et « 2 » sont dans la marge de gauche. Le
feuillet s'arrête sur « parallèle », sans ponctuation ; son verso, vue 143,
est blanc.}

\page{144}
\section{Acta academiæ scientiarum imperialis Petropolitanæ pro anno 1779 Pars prior}

(103)

\emph{Investigatio motuum quibus laminæ et virgæ elasticæ contremiscunt
auctore L. Eulero}

Quanquam hoc argumentum jam pridem tam ab \uncertain{Illustriss.} D.
Bernoulli, quam a me fusius est pertractatum: tamen quia illo tempore neque
principia, unde hujusmodi motus determinari oportet, satis erant exculta,
neque ea analyseos pars, quæ circa functiones binarum variabilium versatur,
satis explorata; actum agere non videbor, si nunc idem argumentum accuratius
investigavero. Praeterea vero etiam tot diversa motuum genera in hujusmodi
corporibus locum habere possunt, quæ accuratiorem enucleationem postulant;
quamobrem hic operam dabo, ut universam hujus rei disquisitionem ex primis
principiis deducam, et clarius, quam quidem ante est factum, proponam.
Imprimis autem omnia diversa motuum genera, quæ quidem occurrere possunt,
dilucide sum expositurus. Quo igitur omnia fiant magis perspicua, duo
præmittam lemmata quorum altero statum æquilibrii, altero vero motus
virgarum utcunque elasticarum et a potentiis quibuscunque \uncertain{que}
sollicitatarum definietur;
\note{copie, de la même main que les feuillets français, d'un mémoire latin
dont le feuillet donne lui-même le recueil, l'année, le titre et l'auteur.
Le titre du recueil est souligné d'un trait court. Le nombre 71 est à l'encre
en haut à droite ; « (103) », entre parenthèses et d'une écriture plus
petite, est au-dessus du titre du mémoire, et semble reporter la page de
l'imprimé copié — c'est une inférence. Le mot lu « Illustriss. » est une
abréviation dont l'initiale ressemble à un « M ». Entre « quibuscunque » et
« que » se tient un trait vertical appuyé, que la lecture ne résout pas.}

\page{145}
(104)

ubi quidem \struck{tam} virgam quam potentias perpetuo in eodem plano sitas
\struck{esse} esse assumo. Demonstrationem autem horum lemmatum non addo,
quoniam eam alio loco jam dedi, atque adeo etiam \uncertain{a} eos casus,
quibus motus non fit in eodem plano, accommodavi.

\subsection{Lemma 1}

§ 1. Si virgæ utcunque elasticæ $EYF$ in singulis elementis potentiæ
quæcunque fuerint applicatæ, statum æquilibrii definire
\marginal{Tab. II fig. 3}

\emph{Solutio.}

Referatur virga ad axem fixum $OV$, et pro ejus puncto quocunque $Y$
statuantur coordinatae orthogonales $OX = x$ $XY = y$, portio autem virgae
$EY = s$, ut sit $ds^2 = dx^2 + dy^2$; tum vero elementi $Yy = ds$ sit
massa $= Sds$, ac in eodem loco elasticitas absoluta $= V$, \struck{ista}
\marginal{ut, posito radio osculi $= r$, vis, seu potius momentum}
elasticitatis sit $\frac{V}{r}$; hincque, loco $r$ substituta formula
differentiali, istud momentum erit
\[
= \frac{V(dy\,ddx - dx\,ddy)}{ds^3},
\]
ubi scilicet nullum differentiale pro constanti est assumtum. Deinde omnes
potentiae, quibus hoc elementum $Yy$ sollicitatur, resolvantur secundum
directiones coordinatarum, ac ponatur vis inde secundum directionem $YP$
resultans $= Pds$ et secundum directionem $YQ = Qds$: quibus positis pro
statu aequilibrii requiritur, ut sit
\note{« (104) » est en haut, au milieu. Les mots « tam » et « esse », en fin
de première ligne, sont barrés d'un trait. « Lemma 1 » et « Solutio » sont
soulignés. L'indication « Tab. II fig. 3 » est dans la marge de gauche, en
regard de l'énoncé du § 1 ; aucune figure n'accompagne la copie. Le mot
« ista » est barré et porte un petit « + » en exposant, qui renvoie à
l'addition portée dans la marge de gauche et précédée du même signe, placée
ici à l'endroit du renvoi. Le mot lu « a » devant « eos casus » est d'une
seule lettre sur la vue.}

\page{146}
(105)
\[
\int dy \int Pds - \int dx \int Qds = V\left(\frac{dy\,ddx - dx\,ddy}{ds^3}\right).
\]
Praeterea vero tensio, qua elementum $Yy$ secundum tangentem
antrorsum versus $E$ tenditur, erit,
$-\left(\frac{dx}{ds}\right)\int Pds - \left(\frac{dy}{ds}\right)\int Qds$.

\emph{Scholion.}

§. 2. Hic assumsimus virgam in statu naturali in directum esse extensam,
ita ut in hoc statu radius osculi ubique sit infinitus. Sin autem virga
naturaliter jam fuerit incurvata, et pro ejus puncto $Y$ radius osculi
ponatur $= \varsigma$, tum, quia vis elasticitatis eatenus tantum se exerit
quatenus ista curvatura in statu naturali sive augetur sive diminuitur, pro
statu æquilibrii habebitur ista aequatio:
$\int dy \int Pds - \int dx \int Qds = V\left(\frac{dy\,ddx - dx\,ddy}{ds^3}\right) - \frac{V}{\varsigma}$.

Caeterum quia formulae integrales $\int Pds$ et $\int Qds$ denotant summam
omnium virium elementarium, portioni virgae $EY = s$ applicatarum,
manifestum est, si virga in ipso termino $E$ a viribus finitis
sollicitetur, eas simul in his formulis integralibus comprehendi debere.

\subsection{Lemma 2}

§. 3 Si eadem virga elastica, quam descripsimus quomodocunque super eodem
plano fuerit projecta ejus motum investigare, hoc est, ejus situm et
figuram ad quodvis tempus definire.

\emph{Solutio}

Maneant omnes denominationes, ut modo sunt constitutae, atque elapso
tempore $= t$
\note{« (105) » en haut, au milieu ; le nombre 72 est à l'encre en haut à
droite. Dans la première formule, le $x$ de « $\int dx$ » est empâté d'une
tache d'encre. Un trait vertical se tient dans la marge devant
« Praeterea ». Le rayon osculateur
de la verge naturellement courbe est noté d'une lettre grecque bouclée, rendue
ici $\varsigma$ ; la même lettre revient aux vues suivantes. Dans les deux
formules du feuillet, la parenthèse ouverte après $V$ embrasse la fraction
entière. « Scholion », « Lemma 2 » et « Solutio » sont soulignés. La ligne
s'arrête sur « tempore = t », que la vue 147 poursuit.}

\page{147}
(106)

(perpetuo in minutis secundis exprimendo) teneat virga figuram in tabula
repraesentatam, ita ut portioni $EY = s$ respondeant coordinatae $OX = x$ et
$XY = y$, quae, quia cum tempore continuo variantur, hic tanquam functiones
duarum variabilium $s$ et $t$ spectari debent. Hinc igitur \emph{colligantur}
formulae $\left(\frac{ddx}{dt^2}\right)$ et $\left(\frac{ddy}{dt^2}\right)$
quibus uncinulis $(\;)$ indicatur, in utraque differentiatione solum tempus
$t$ variabile esse assumtum, arcum vero $s$ pro constante esse habitum. Hinc
jam ex viribus elementaribus, virgae in singulis punctis applicatis,
formentur sequentes valores:
\[
P' = P - \frac{S}{2g}\left(\frac{ddx}{dt^2}\right) \quad\text{et}\quad
Q' = Q - \frac{S}{2g}\left(\frac{ddy}{dt^2}\right),
\]
ubi $g$ denotat altitudinem lapsus gravium in uno minuto secundo: atque ex
his status virgae pro hoc tempore isthac exprimetur aequatione:
\[
\int dy \int P'ds - \int dx \int Q'ds
= V\left(\frac{dy\,ddx - dx\,ddy}{ds^3}\right) - \frac{V}{\varsigma},
\]
si quidem virga in statu naturali jam ita fuerit incurvata, ut ejus puncto
$Y$ conveniat radius osculi $\varsigma$; unde si virga fuerit naturaliter
recta, ob $\varsigma = \infty$ terminus postremus $\frac{V}{\varsigma}$
omitti poterit. Denique, quod ad tensionem in singulis punctis attinet, erit
simili modo quo supra tensio in $Y$ versus $E$ vergens
\[
-\left(\frac{dx}{ds}\right)\int P'ds - \left(\frac{dy}{ds}\right)\int Q'ds.
\]
his igitur formulis, siquidem eas evolvere licuerit, totus virgae motus
determinabitur.

\subsection{Problema 1}

§. 4 Si virgae datae longitudinis $EF$, naturaliter recta et ubique tam
\marginal{Tab. II fig. 4}
\note{« (106) » en haut, au milieu. Un trait vertical sépare « minutis » de
« secundis » à la première ligne. Le $O$ de « $OX = x$ » est empâté d'une
tache. « colligantur » est souligné sur le feuillet. Dans la première
intégrale de l'équation, le $P'$ est écrit sur une lettre reprise. Le
dénominateur lu « $2g$ » est écrit d'un chiffre bouclé. « Problema 1 » est
souligné ; l'indication « Tab. II fig. 4 » est dans la marge de gauche, en
regard du § 4.}

\page{148}
(107)

aequaliter crassa quam aequaliter elastica, utcunque contremiscat,
aequationem invenire, qua omnes motus, qui in virga locum habere possunt
continentur

\emph{Solutio.}

Referat igitur $EF$ virgam nostram in statu naturali constitutam, cujus
longitudo sit $EF = a$, ejusque crassities $= cc$, ita ut ejus volumen sit
$acc$; ac per talia volumina tam massas quam vires sollicitantes exprimamus,
ita ut si dicamus quampiam vim esse $= b^3$ ea tanta sit intelligenda,
quantum foret pondus ejusdem materiae, ex qua virga constat, sub volumine
$b^3$ contentum. Hinc si in virga capiatur portio $EX = x = s$, quandoquidem
in vibrationibus minimis arcus $s$ perpetuo abscissae $x$ aequari potest,
elementi $Xx = ds$ massula erit $cc\,ds$, ita ut, quod supra vocavimus $S$,
hic nobis sit $cc$. Tempore autem elapso $t$, ob tremorem conceptum
transierit punctum $X$ in $Y$, posito $XY = y$. et quia vibrationes quam
minimae statuantur, ista applicata $y$ prae abscissa $EX = x$, sive arcu
$s$, quasi evanescit; neque idem punctum $x$ alium motum habere nequit, nisi
in directione applicatae $XY$; unde motus secundum directionem $x$ erit
nullus, ideoque $\frac{ddx}{dt^2} = 0$; atque hinc ob $ddx = 0$, radius
osculi erit
\[
\frac{ds^3}{-dx\,ddy} = -\frac{ds^2}{ddy};
\]
praeterea vero erit $\varsigma = \infty$. Quod autem ad elasticitatem
absolutam attinet, ea pro similibus virgis recte quadrato crassitiei
reputatur proportionalis, ita ut sit $V$ \uncertain{at} $c^4$, unde
statuabimus $V = bc^4$ ubi $b$ denotat quantitatem ab indole
\note{« (107) » en haut, au milieu ; le nombre 73 est à l'encre en haut à
droite. Un trait vertical se tient entre « elastica » et « utcunque ».
« Solutio » est souligné. Dans la formule du rayon osculateur, l'exposant de
$ds$ au numérateur est écrit très haut au-dessus de la lettre, et lu 3. Entre
$V$ et $c^4$, à l'avant-dernière ligne, le feuillet porte deux lettres, lues
« at ». Deux taches d'encre sont dans la marge de gauche, en regard de
« praeterea ».}

\page{149}
(108)

materiae virgae pendentem, et quia $\frac{V}{r}$ momentum virium exhibet,
vis autem nobis per volumen denotatur, formula $\frac{V}{r}$ quatuor
dimensiones lineares complecti debet; unde patet, literam \emph{b} certam
longitudinem referre.

Quia porro virgam a nullis viribus elementaribus urgeri statuimus, erit tam
$P = 0$ quam $Q = 0$. interim tamen, si sumamus virgam in termino $E$ duas
sustinere vires, alteram in directione $EE = E$, alteram in directione
$EF = E$, alteram vero in directione $Ef = E$, fieri debebit
\[
\int Pds = E \quad\text{et}\quad \int Qds = F.
\]
Cum igitur ob $S = cc$ poni oporteat
\[
P' = P - \frac{cc}{2g}\left(\frac{ddx}{dt^2}\right) = F, \quad
\text{ob}\ \left(\frac{ddy}{dt^2}\right) = 0,
\]
sit $\int P'ds = E$; tum vero erit
\[
Q' = Q - \frac{cc}{2g}\left(\frac{ddy}{dt^2}\right), \ \text{hincque porro}
\]
\[
\int Q'ds = F - \frac{cc}{2g}\int ds \left(\frac{ddy}{dt^2}\right).
\]
His igitur valoribus substitutis aequatio pro motu virgae induet hanc
formam:
\[
Ey - Fx + \frac{cc}{2g}\int dx \int ds \left(\frac{ddy}{dt^2}\right)
= -\frac{bc^4\,ddy}{ds^2}.
\]
Pro tensione autem, qua punctum $Y$ versus $E$ tenditur, habebimus
$-E - \left(\frac{dy}{ds}\right)\int Q'ds$, ubi, quia $\frac{dy}{ds}$ quasi
evanescit, tensio simpliciter erit $-E$, unde casu quo $E = 0$ tensio ubique
etiam erit nulla.

Differentiemus igitur aequationem pro motu inventam, sumto solo elemento
$ds = dx$ constante, fietque
\[
E\,dy - F\,dx + \frac{cc}{2g}\,dx \int ds \left(\frac{ddy}{dt^2}\right)
= -\frac{bc^4\,d^3y}{ds^2}.
\]
et per $ds$ dividendo
\[
\frac{E\,dy}{ds} - F + \frac{cc}{2g}\int ds \left(\frac{ddy}{dt^2}\right)
= -\frac{bc^4\,d^3y}{ds^3};
\]
\note{« (108) » en haut, au milieu. Un trait vertical se tient entre
« virium » et « exhibet » ; la lettre \emph{b} de « literam b » est
soulignée. La phrase sur les forces du terme $E$ porte trois membres
« alteram » pour « duas vires », avec les directions lues $EE$, $EF$ et
$Ef$, chacune « $= E$ » : les lettres ont été lues sur la vue et ne sont
pas corrigées. Dans la valeur de $P'$, le membre après le second signe
d'égalité est lu $F$, et la parenthèse qui suit « ob » est lue $ddy$ ; la
tête de ces lettres est serrée contre la ligne précédente. Dans
« E dy − F dx », la seconde majuscule se distingue mal d'un $E$. La
terminaison de « elemento » est écrite sur une autre, reprise à la plume.
La dernière formule finit par un point-virgule, tracé sous la ligne.}

\page{150}
(109)

haec aequatio denuo differentiata ac per $ds$ divisa praebebit istam:
\[
E\frac{ddy}{ds^2} + \frac{cc}{2g}\left(\frac{ddy}{dt^2}\right) = -\frac{bc^4\,d^4y}{ds^4},
\]
quae, quo discrimen inter binas variabiles $s$ et $t$ clarius ob oculos
ponatur, more solito ita repraesentetur:
$E\left(\frac{ddy}{ds^2}\right) + \frac{cc}{2g}\left(\frac{ddy}{dt^2}\right) = -bc^4\left(\frac{d^4y}{ds^4}\right)$.

\subsection{Corollarium. I}

§. 5 Quod si ergo virga a nullis plane viribus externis extendatur, ita ut
sit $E = 0$, aequatio motum virgae determinans erit
\[
\frac{cc}{2g}\left(\frac{ddy}{dt^2}\right) = -bc^4\left(\frac{d^4y}{ds^4}\right), \ \text{sive}
\]
\[
\frac{1}{2g}\left(\frac{ddy}{dt^2}\right) = -bc^2\left(\frac{d^4y}{ds^4}\right),
\]
ita ut totum negotium huc sit reductum, quemadmodum ista aequatio
differentialis quarti gradus integrari queat; ubi quidem in limine
confiteri cogimur, ejus integrale nullo adhuc modo inveniri potuisse, ita ut
contenti esse debeamus in solutiones particulares inquirere.

\subsection{Corollarium. II}

§. 6. Sin autem accedat vis litera $E$ indicata, ejus duo casus perpendendi
occurrunt: \uncertain{prouti} virga ab ea vel extenditur vel comprimitur.
Talis enim virga quatenus est rigida, etiam vires sustinere potest, quae
ipsam comprimere conantur, cujusmodi vires eo majores esse possunt quo major
fuerit elasticitas; si enim esset perfecte flexilis, nulla plane vis
comprimens admitti posset, unde pro quovis elasticitatis gradu indagandum
erit, quantam vim comprimentem sustinere valeat, antequam incurvetur
\note{« (109) » en haut, au milieu ; le nombre 74 est à l'encre en haut à
droite. Un trait vertical se tient dans la marge devant « haec ».
« Corollarium. I » et « Corollarium. II » sont soulignés. Dans la seconde
forme de l'équation du § 5, le coefficient est lu $bc^2$, là où la première
porte $bc^4$. Le mot lu « prouti » est écrit avec deux $t$ apparents. Un
petit signe à l'encre est dans la marge de gauche, en regard de
« possunt ».}

\page{151}
(110)

quam quidem quaestionem jam olim solutam dedi, ubi vires, quas columnae
sustentare valeant, sum perscrutatus. Quod autem ad alteram vim
extendentem attinet, ab ea virga, quasi esset chorda perfecte flexilis,
extendi poterit. Concipiatur scilicet ejus termino $E$ filum seu chorda
alligata, quae per foraminulum $O$ fulcri immobilis $mn$ traducta circa
trochleam $T$ certum pondus $P$ habeat appensum. Hoc igitur modo virga non
solum ob propriam elasticitatem, sed etiam ob pondus tendens $P$ motum
tremulum concipiet, unde sonum edet mixtum seu medium quendam inter sonum
virgae elasticae proprium et sonum chordæ tensae.
\marginal{tab. II fig. 5}

\subsection{Corollarium III}

§. 7 Ponamus igitur hujusmodi vim tendentem esse $= cch$, et quia in plagam
contrariam dirigitur, erit $E = -cch$, unde pro isto motu habebimus
sequentem aequationem:
\[
-h\left(\frac{ddy}{ds^2}\right) + \frac{1}{2g}\left(\frac{ddy}{dt^2}\right)
= -bcc\left(\frac{d^4y}{ds^4}\right),
\]
quae ergo aequatio, si fuerit $b = 0$, quo casu elasticitas evanescit,
motum chordae perfecte flexilis exprimet; erit enim
\[
\frac{1}{2g}\left(\frac{ddy}{dt^2}\right) = h\left(\frac{ddy}{ds^2}\right)
\]
quemadmodum contra si $h = 0$, orietur sonus virgae elasticae proprius,
fietque
$\frac{1}{2g}\left(\frac{ddy}{dt^2}\right) = -bcc\left(\frac{d^4y}{ds^4}\right)$
et si casus ex utroque fuerit mixtus, habebimus
$\frac{1}{2g}\left(\frac{ddy}{dt^2}\right) = h\left(\frac{ddy}{ds^2}\right) - bcc\left(\frac{d^4y}{ds^4}\right)$
\note{« (110) » en haut, au milieu. L'indication « tab. II fig. 5 » est dans
la marge de gauche, en regard de « flexilis, extendi poterit » ; la lettre
$O$ de « foraminulum O » et les lettres $mn$ sont soulignées. Un trait
d'encre est dans la marge, en regard de la première ligne. « Corollarium
III » est souligné. Les dernières formules sont écrites dans la ligne, à
droite du texte, sans ponctuation finale.}

\page{152}
(111)

\subsection{Scholion}

§. 8. Quoniam igitur vires secundum ipsam virgae directionem agentes,
quibus ea vel extenderetur vel comprimeretur, ad aequationem finalem magis
complicatam perducunt, eas in hac tractatione penitus removeamus, et nostras
investigationes ad eos tantum motus restringamus, quos virga, sive a nullis
viribus sollicitata, sive a talibus tantum, quae in virgae directionem
normaliter agunt, quas supra litera $F$ designavimus, recipere potest. Ipsam
vero virgam hic perpetuo naturaliter rectam et per totam longitudinem ubique
aequaliter crassam et aequaliter elasticam statuamus, easdem denominationes
retinentes quae hactenus sunt descriptae. At quia aequatio finalis per
duplicem differentiationem est orta, etiam praecedentes aequationes, quae ad
eam deduxerunt considerasse juvabit, quae sunt:
\[
\text{I.} \quad -Fx + \frac{cc}{2g}\int dx \int ds \left(\frac{ddy}{dt^2}\right) = -bc^4\left(\frac{ddy}{ds^2}\right)
\]
\[
\text{II.} \quad -F + \frac{cc}{2g}\int dx \left(\frac{ddy}{dt^2}\right) = -bc^4\left(\frac{d^3y}{ds^3}\right)
\]
\[
\text{III.} \quad \frac{cc}{2g}\left(\frac{ddy}{dt^2}\right) = -bc^4\left(\frac{d^4y}{ds^4}\right), \ \text{sive}
\]
\[
\frac{1}{2g}\left(\frac{ddy}{dt^2}\right) = -bcc\left(\frac{d^4y}{ds^4}\right)
\]
ubi in priores adhuc vis $F$, qua virga normaliter urgeri potest,
ingreditur, cujus ratio erit habenda, quando virga non omnino est libera, sed
in uno pluribusve punctis quasi ope styli est fixa, circa quem tamen sit
mobilis; unde statim intelligitur, virgam infinitis modis per tales stylos
affigi posse, quibus ejus motus diversimode temperetur; quos diversos casus in
sequentibus accuratius evolvemus.
\note{« (111) » en haut, au milieu, et au-dessous « Scholion », souligné,
précédé d'un petit trait vertical ; le nombre 75 est à l'encre en haut à
droite. Les trois équations sont numérotées I, II, III dans la ligne, en
chiffres romains.}

\page{153}
(112)

\subsection{Problema}

§. 9. Cum virga nostra infinitis modis contremiscere queat, eos casus in
genere investigare, quibus ejus motus vibratorius evadet regularis, seu
minimis oscillationibus penduli cujuspiam simplicis conformis

\emph{Solutio.}

Ponamus igitur longitudinem istius penduli simplicis $= k$, atque ut motus
virgae pari modo peragatur, necesse est ut sit
$\frac{1}{2g}\left(\frac{ddy}{dt^2}\right) = -\frac{y}{k}$: in qua aequatione
cum solum tempus $t$ pro variabili habeatur, longitudo vero \emph{s} ut
constans spectetur, per duplicem integrationem pervenitur ad istam formulam
generalem:
\[
y = \mathfrak{M}\,\text{sin.}\left(\zeta + t\frac{\sqrt{2g}}{k}\right),
\]
ubi quantitas $\mathfrak{M}$ non solum est constans, sed etiam functionem
quamcunque ipsius \emph{s} designare potest. Cum igitur per aequationem
finalem sit
\[
\frac{1}{2g}\left(\frac{ddy}{dt^2}\right) = -bcc\left(\frac{d^4y}{ds^4}\right),
\]
erit etiam $bcc\left(\frac{d^4y}{ds^4}\right) = \frac{y}{k}$: in qua equatione
sola quantitas \emph{s} pro variabili assumitur, cujus ergo integrale
investigari oportet. Quod quo facilius fieri possit ponamus brevitatis gratia
$bcck = f^4$, et remotis \uncertain{clausulis} jam sit $y = \frac{f^4\,d^4y}{ds^4}$, cui
satisfieri posse evidens est hujusmodi valore: $y = e^{\lambda s}$. Quia enim
hinc fit
\[
\frac{dy}{ds} = \lambda e^{\lambda s}, \quad
\frac{ddy}{ds^2} = \lambda\lambda e^{\lambda s}, \quad
\frac{d^3y}{ds^3} = \lambda^3 e^{\lambda s}, \quad
\frac{d^4y}{ds^4} = \lambda^4 e^{\lambda s},
\]
his substitutis prodit ista aequalitas: $1 = \lambda^4 f^4$ sive
$\lambda^4 = \frac{1}{f^4}$
\note{« (112) » en haut, au milieu. « Problema » et « Solutio » sont
soulignés ; les lettres \emph{s} de « longitudo vero s », « ipsius s » et
« sola quantitas s » sont soulignées sur le feuillet, et rendues ici en
italique. Le coefficient de la formule générale est une majuscule
calligraphiée, proche d'un « M » ouvert, rendue $\mathfrak{M}$ ; la lettre
de l'exposant, rendue $\lambda$, est une petite lettre à jambage croisé. Dans
$t\frac{\sqrt{2g}}{k}$, le radical ne couvre que $2g$, et $k$ est écrit sous
la barre. « remotis clausulis » est lu sur un mot serré. Le nombre 1 de
« $1 = \lambda^4 f^4$ » et la fin de la dernière ligne touchent le bord
inférieur du feuillet.}

\page{154}
(113)

unde sequentes quatuor valores pro litera $\lambda$ elicientur, scilicet:
\[
1^\circ.\ \lambda = \frac{1}{f}; \quad 2^\circ.\ \lambda = -\frac{1}{f}; \quad
3^\circ.\ \lambda = \frac{\sqrt{-1}}{f}; \quad 4^\circ.\ \lambda = -\frac{\sqrt{-1}}{f};
\]
quibus valoribus inventis constat, eos etiam utcunque combinatos satisfacere,
ita ut statuere queamus
\[
y = \alpha e^{\frac{s}{f}} + \beta e^{-\frac{s}{f}} + \gamma e^{\frac{s\sqrt{-1}}{f}} + \delta e^{-\frac{s\sqrt{-1}}{f}};
\]
quae forma, cum quatuor contineat constantes arbitrarias, utique continet
integrale completum nostrae aequationis. Notum autem est bina posteriora
membra imaginaria reduci ad sinum et cosinum anguli realis $\frac{s}{f}$. Hinc
igitur multiplicando per communem factorem $N$, qui, ob tempus \emph{t}
constans assumtum, pro functione temporis quacunque haberi potest, habebimus
\[
y = N\left(\alpha e^{\frac{s}{f}} + \beta e^{-\frac{s}{f}} + \gamma\,\text{sin.}\frac{s}{f} + \delta\cos.\frac{s}{f}\right);
\]
qui ergo valor ante invento
$y = \mathfrak{M}\,\text{sin.}\left(\zeta + t\frac{\sqrt{2g}}{k}\right)$
aequalis reddi debet, id quod facillime praestabitur statuendo:
\[
N = C\,\text{sin}\left(\zeta + t\frac{\sqrt{2g}}{k}\right) \ \text{et}
\]
\[
\mathfrak{M} = C\left(\alpha e^{\frac{s}{f}} + \beta e^{-\frac{s}{f}} + \gamma\,\text{sin.}\frac{s}{f} + \delta\cos\frac{s}{f}\right),
\]
sic enim uterque valor reducetur ad istam aequationem:
\[
y = C\,\text{sin}\left(\zeta + t\frac{\sqrt{2g}}{k}\right)\left(\alpha e^{\frac{s}{f}} + \beta e^{-\frac{s}{f}} + \gamma\,\text{sin.}\frac{s}{f} + \delta\cos.\frac{s}{f}\right);
\]
ubi jam literae $C, \alpha, \beta, \gamma, \delta$ cum angulo $\zeta$ denotant
veras quantitates constantes pro arbitrio accipiendas; quomodocunque autem
accipiantur motus virgae semper erit comparatus, ut congruat oscillationibus
penduli
\note{« (113) » en haut, au milieu ; le nombre 76 est à l'encre en haut à
droite. Les lettres grecques sont d'une main peu appuyée : $\alpha$ ressemble
à un « u » et $\gamma$ à un « v » ou à un $\vartheta$ ; la lecture les tient
pour les mêmes lettres d'un bout à l'autre. La lettre \emph{t} de « tempus
t » est soulignée. La majuscule du second membre de l'égalité
« $\mathfrak{M} = C(\dots)$ » est de la même forme que le coefficient de la
formule générale de la vue 153.}

\page{155}
(114)

simplicis, cujus longitudo est $= k$, unde tempus cujusque vibrationis erit
$= \pi\sqrt{\frac{k}{2g}}$, expressum in minutis secundis, hincque porro
numerus vibrationum uno minuto secundo editarum erit
$= \frac{\sqrt{2g}}{\pi\sqrt{k}}$, qui numerus etiam pro mensura soni, quem
chorda edet, haberi solet.

\subsection{Corollarium 1}

§. 10. formulae illae exponentiales, perinde ac sinus et cosinus, commode per
series infinitas exhiberi possunt, quae ita se habebunt:
\[
y = C\,\text{sin.}\left(\zeta + t\frac{\sqrt{2g}}{k}\right)
\left\{
\begin{array}{l}
\alpha(1 + \frac{s}{f} + \frac{ss}{1.2ff} + \frac{s^3}{1.2.3f^3} + \frac{s^4}{1.2.3.4f^4} + \&c. \\[4pt]
+\beta(1 - \frac{s}{f} + \frac{ss}{1.2ff} - \frac{s^3}{1.2.3.f^3} + \frac{s^4}{1.2.3.4f^4} - \&c. \\[4pt]
+\gamma(\phantom{1 +{}} \frac{s}{f} \qquad\qquad - \frac{s^3}{1.2.3} \qquad\qquad + \&c \\[4pt]
+\delta(1 \qquad\quad - \frac{s^2}{1.2.ff} \qquad\qquad + \frac{s^4}{1.2.3.4.f^4}
\end{array}
\right.
\]
Hinc igitur, si brevitatis gratia ponamus
\[
y = C\,\text{sin.}\left(\zeta + t\frac{\sqrt{2g}}{k}\right)S, \ \text{erit per seriem infinitam}
\]
\[
S = (\alpha + \beta + \delta)1 + (\alpha - \beta + \gamma)\frac{s}{f} + (\alpha + \beta - \delta)\frac{ss}{1.2ff}
\]
\[
+ (\alpha - \beta - \gamma)\frac{s^3}{1.2.3.f^3} + (\alpha + \beta + \delta)\frac{s^4}{1.2.3.4f^4} + (\alpha - \beta + \gamma)\frac{s^5}{1.2.3.4.5f^5} + \&c.
\]
Hae igitur series eo magis convergent, quo major fuerit quantitas \emph{f}
prae arcu \emph{s}; ubi recordemur esse $f = \sqrt[4]{bcck}$, ita ut simul
longitudinem penduli simplicis \emph{k} in se complectatur.

\subsection{Corollarium. 2.}

§ 11 Invento valore ipsius $y$, ejus quoque valores per differentiationem
\note{« (114) » en haut, au milieu. Les quatre séries sont réunies par une
grande accolade ; les parenthèses ouvertes après $\alpha$, $\beta$, $\gamma$,
$\delta$ ne sont pas refermées, et les termes sont disposés en colonnes, que
la transcription imite. Au terme en $s^3$ de la ligne de $\gamma$, le
dénominateur est écrit « 1.2.3 » sans $f^3$ ; la ligne de $\delta$ ne porte
pas de « \&c ». Les signes finals, rendus « \&c. », ont une forme bouclée
proche de $\infty$. Les lettres \emph{f}, \emph{s} et \emph{k} de la
dernière phrase sont soulignées sur le feuillet. « Corollarium 1 » et
« Corollarium. 2. » sont soulignés.}

\page{156}
(115)

eruti assignari possunt: reperietur autem
\[
\left(\frac{dy}{ds}\right) = \frac{C}{f}.\,\text{sin.}\left(\zeta + t\frac{\sqrt{2g}}{k}\right)\left(\alpha e^{\frac{s}{f}} - \beta e^{-\frac{s}{f}} + \gamma\cos.\frac{s}{f} - \delta\,\text{sin.}\frac{s}{f}\right)
\]
\[
\left(\frac{ddy}{ds^2}\right) = \frac{C}{ff}.\,\text{sin.}\left(\zeta + t\frac{\sqrt{2g}}{k}\right)\left(\alpha e^{\frac{s}{f}} + \beta e^{-\frac{s}{f}} - \gamma\,\text{sin.}\frac{s}{f} - \delta\cos.\frac{s}{f}\right)
\]
\[
\left(\frac{d^3y}{ds^3}\right) = \frac{C}{f^3}.\,\text{sin}\left(\zeta + t\frac{\sqrt{2g}}{k}\right)\left(\alpha e^{\frac{s}{f}} - \beta e^{-\frac{s}{f}} - \gamma\cos.\frac{s}{f} + \delta\,\text{sin.}\frac{s}{f}\right)
\]
\[
\left(\frac{d^4y}{ds^4}\right) = \frac{C}{f^4}.\,\text{sin.}\left(\zeta + t\frac{\sqrt{2g}}{k}\right)\left(\alpha e^{\frac{s}{f}} + \beta e^{-\frac{s}{f}} + \gamma\,\text{sin.}\frac{s}{f} + \delta\cos.\frac{s}{f}\right)
\]

\subsection{Corollarium. 3.}

§. 12. Quod si autem tantum tempus \emph{t} variabile accipiamus, ipsum
motum cognoscemus, quo singulae virgae ciebuntur; namque celeritas, qua
virgae punctum $X$ in directione $XY$ movetur, est
$\left(\frac{dy}{dt}\right)$, quae ergo per valorem inventum erit
\[
\left(\frac{dy}{dt}\right) = \frac{C\sqrt{2g}}{\sqrt{k}}\cos\left(\zeta + t\frac{\sqrt{2g}}{k}\right)\left(\alpha e^{\frac{s}{f}} + \beta e^{-\frac{s}{f}} + \gamma\,\text{sin.}\frac{s}{f} + \delta\cos.\frac{s}{f}\right)
\]
et quia incrementum celeritatis, per elementum temporis \emph{dt} divisum,
praebet accelerationem, erit ea
\[
\left(\frac{ddy}{dt^2}\right) = -\frac{2Cg}{k}\,\text{sin}\left(\zeta + t\frac{\sqrt{2g}}{k}\right)\left(\alpha e^{\frac{s}{f}} + \beta e^{-\frac{s}{f}} + \gamma\,\text{sin.}\frac{s}{f} + \delta\cos.\frac{s}{f}\right);
\]
ex quibus formulis perspicitur, quomodo omnibus conditionibus praescriptis
satisfiat.

\subsection{Scholion I}

§. 13 Quanquam igitur hic tantum motus vibratorios regulares contemplamur:
tamen ob ingentem quantitatum constantium arbitrariarum numerum infinita
varietas locum habere potest, prouti istae constantes aliter atque aliter
determinantur. In hoc autem negotio imprimis
\note{« (115) » en haut, au milieu, le premier chiffre mal formé ; le nombre
77 est à l'encre en haut à droite. Un trait vertical se tient entre
« reperietur » et « autem ». Les lettres \emph{t} de « tempus t » et
\emph{dt} de « temporis dt » sont soulignées. La fin de « perspicitur » et
une lettre de « determinantur » sont empâtées d'une tache. « Corollarium.
3. » et « Scholion I » sont soulignés.}

\page{157}
(116)

ad statum virgae est respiciendum, utrum ea perfecte sit libera et nullis
plane viribus coerceatur, an vero alicubi sit fixa vel ope unius pluriumve
stylorum. Statim enim ac virgae status fuerit definitus, etiam constantium
$\alpha, \beta, \gamma, \delta$ ratio non solum perfecte determinatur, sed
etiam obtinebitur aequatio, ex qua valorem quantitatis \emph{f} determinare
licebit, hincque igitur ipsa penduli simplicis longitudo $k$ elicietur.
Praecipue autem in hoc negotio ad totam virgae longitudinem $EF = a$, cujus
hactenus nulla ratio est habita, respici oportet. Quod si autem nullas alias
vires $F$ admittamus, nisi quae virgae in utroque termino sint applicatae,
status utriusque termini triplex occurrere potest, quos igitur casus hic
evolvere conveniet; ubi quod de termino $E$ definiemus simul pro altero $F$
intelligi debet.

I. Primus igitur casus esto, quo virgae terminus $E$ plane est liber et a
nulla vi coercetur; tum igitur posito $s = 0$, tribus aequationibus supra
(§. 8) memoratis satisfieri oportet, unde, quia vis $F = 0$ et utrumque
integrale
\[
\int ds \left(\frac{ddy}{dt^2}\right) \quad\text{et}\quad \int ds \int ds \left(\frac{ddy}{dt^2}\right)
\]
perpetuo ita accipi supponitur, ut evanescat posito $s = 0$, ex prima pro
termino $E$, ubi $s = 0$, oportet esse $\left(\frac{ddy}{ds^2}\right) = 0$; ex
secunda autem equatione nascitur ista conditio:
$\left(\frac{d^3y}{ds^3}\right) = 0$, ita ut hic duas determinationes
\uncertain{portabit}.

II. Secundus casus esto, quo virgae terminus $E$ ope styli ita figitur, ut
circa eum libere moveri possit: quo ergo casu vis quaedam $F$ stilum retinens
\note{« (116) » en haut, au milieu. Le trait qui prolonge « unius » jusqu'au
bout de la ligne est un trait de remplissage. Dans « quia vis F = 0 », la
majuscule se distingue mal d'un $E$ ; elle est lue $F$, comme la force
normale du § 8. Dans les deux intégrales, le dénominateur est écrit serré et
se lit mal : $dt^2$ est la lecture la plus probable. Dans la seconde
condition, la lettre du numérateur de $\frac{d^3y}{ds^3}$ ressemble à un
« s ». Le dernier mot de l'avant-dernier alinéa, lu « portabit », touche le
bord droit du feuillet et finit par un point. Le chiffre du « II » est repris
à la plume sur un « I ».}

\page{158}
(117)

aderit. Primum igitur, quia terminus $E$ in suo loco fixus detinetur, posito
$s = 0$, hoc casu fieri debet $y = 0$; praeterea vero prima memoratarum
aequationum suppeditat hanc conditionem: $\left(\frac{ddy}{ds^2}\right) = 0$;
secunda aequatio determinabit ipsam vim
\[
F = +bc^4\left(\frac{d^3y}{ds^3}\right);
\]
quam autem nosse parum nobis refert; quia in motum vibratorium non influit;
sicque iste casus etiam duas continet determinationes, scilicet:
\[
y = 0 \quad\text{et}\quad \left(\frac{ddy}{ds^2}\right) = 0;
\]
hunc igitur vocemus simpliciter fixum.

III. Tertius casus esto, quo terminus virgae $E$ ita muro quasi est
infixus, ut hoc loco non moveri sed tantum incurvari queat. Hoc igitur casu
manifestum est, posito $s = 0$ non solum fieri debere $y = 0$, sed etiam
$\left(\frac{dy}{ds}\right) = 0$, quia in hoc termino tangens in ipsum axem
incidere debeat; hunc autem casum vocemus infixum.

Cum igitur hi singuli casus binas determinationes involvant, si etiam ad
alterum terminum $F$ respiciamus, pro quo habebimus $x = s = a$, quicunque
casus in utroque locum habeant, semper inde quatuor oriuntur determinationes,
quibus non solum constantes illae $\alpha, \beta, \gamma, \delta$ definientur,
sed insuper aequatio resultabit, ex qua ipsam quantitatem \emph{f}, hincque
pendulum simplex $k$ assignare licebit

\subsection{Scholion 2}

§. 14. Cum igitur pro utroque termino $E$ et $F$ terni casus locum habere
queant, hinc omnino sex casus diversi nascuntur, quos ordine
\uncertain{evolvi} conveniet, siquidem hoc argumentum in omni
\uncertain{extensione} pertractare voluerimus; hos igitur sex casus sequenti
modo designabimus:
\note{« (117) » en haut, au milieu ; le nombre 78 est à l'encre en haut à
droite. « prima memoratarum » est écrit « primat memoratarum », le $t$
barrant la jonction. Le chiffre « sex », deux fois, est écrit d'une
initiale proche d'un « 2 ». « Scholion 2 » est souligné. Le feuillet
s'arrête sur les deux points, et la liste annoncée commence à la vue 159.}

\page{159}
(118)

\begin{itemize}
  \item[I] terminus $E$ liber; \quad terminus $F$ liber
  \item[II.] ---\ $E$ liber; --- $F$ simplic. fixus
  \item[III] ---\ $E$ liber; --- $F$ infixus
  \item[IV] ---\ $E$ simpliciter fixus; --- $F$ simplic. fixus
  \item[V.] ---\ $E$ simpliciter fixus; --- $F$ infixus
  \item[VI.] ---\ $E$ infixus; --- $F$ infixus.
\end{itemize}

Plures quidem casus videntur occurrere, veluti si $E$ esset fixus et $F$
liber; sed quia ambos terminos inter se permutare licet, hic casus in illis
memoratis sex jam continetur. His igitur positis sex nobis supersunt
problemata, quae breviter pertractabimus.

\section{Evolutio casus I quo virgae elasticae uterque terminus prorsus est liber}

\subsection{Problema}

§. 15 Si virga elastica a nullis plane viribus sollicitetur et plano
horizontali politissimo libere incumbat, investigare omnes vibrationes
regulares, quibus ea contremiscere possit.

\emph{Solutio}

Quoniam uterque terminus $E$ et $F$ liber assumitur, pro utroque erit tam
$\left(\frac{ddy}{ds^2}\right) = 0$ quam $\left(\frac{d^3y}{ds^3}\right) = 0$,
hinc ergo pro termino $E$ posito $s = 0$ nascuntur hae duae aequationes:
\[
\text{I.}\ \alpha + \beta - \delta = 0; \qquad \text{II.}\ \alpha - \beta - \gamma = 0
\]
pro altero termino $F$ erit $s = a$, et posito brevitatis gratia
$\frac{a}{f} = \omega$ istae oriuntur aequatio
\note{« (118) » en haut, au milieu. Les six cas sont écrits en tableau, les
chiffres romains dans la marge, chaque ligne reliée à la colonne de droite
par des tirets ; les tirets de début de ligne tiennent lieu de « terminus ».
Les deux lignes du titre « Evolutio casus I / quo virgae elasticae uterque
terminus prorsus est liber » sont centrées, la seconde soulignée d'un bout à
l'autre ; « Problema » et « Solutio » sont soulignés. Dans la seconde
condition, la lettre du numérateur de $\frac{d^3y}{ds^3}$ ressemble de nouveau
à un « s » (voir la vue 157). Un trait vertical suit « hae duae
aequationes: ». Le dernier mot, « aequatio », arrive au bord droit du
feuillet, qui en cache la fin ; la phrase se poursuit à la vue 160.}

\page{160}
(119)
\[
\text{III.}\ \alpha e^{\omega} + \beta e^{-\omega} - \gamma\,\text{sin.}\,\omega - \delta\cos.\omega = 0;
\]
\[
\text{IV.}\ \alpha e^{\omega} - \beta e^{-\omega} - \gamma\cos.\omega + \delta\,\text{sin.}\,\omega = 0.
\]
Jam ex duabus prioribus colligitur $\delta = \alpha + \beta$ et
$\gamma = \alpha - \beta$, qui valores in duabus posterioribus substituti
praebent:
\[
\text{III.}\ \alpha(e^{\omega} - \text{sin.}\,\omega - \cos.\omega) + \beta(e^{-\omega} + \text{sin.}\,\omega - \cos.\omega) = 0;
\]
\[
\text{IV.}\ \alpha(e^{\omega} + \text{sin.}\,\omega - \cos.\omega) - \beta(e^{-\omega} - \text{sin.}\,\omega - \cos.\omega) = 0;
\]
unde duplici modo concluditur:
\[
\frac{\alpha}{\beta} = \frac{-e^{-\omega} - \text{sin.}\,\omega + \cos.\omega}{e^{\omega} - \text{sin.}\,\omega - \cos.\omega}
= \frac{e^{-\omega} - \text{sin.}\,\omega - \cos.\omega}{e^{\omega} + \text{sin.}\,\omega - \cos.\omega}.
\]
Statuamus hic brevitatis gratia $\text{sin.}\,\omega + \cos.\omega = p$ et
$\text{sin.}\,\omega - \cos\omega = q$, ut habeamus
\[
\frac{-e^{-\omega} - q}{e^{\omega} - p} = \frac{e^{-\omega} - p}{e^{\omega} + q};
\]
unde deducitur haec aequatio:
\[
-1 - qq - qe^{\omega} - qe^{-\omega} = 1 + pp - pe^{\omega} - pe^{-\omega}, \ \text{sive}
\]
\[
2 + pp + qq + (p - q)e^{\omega} + (p - q)e^{-\omega} = 0.
\]
Cum igitur sit
\[
pp + qq = 2 \quad\text{et}\quad p - q = -2\cos\omega, \ \text{erit}
\]
\[
2 - \cos\omega(e^{\omega} + e^{-\omega}) = 0 \ \text{hincque}
\]
\[
\cos\omega = \frac{2}{e^{\omega} + e^{-\omega}};
\]
quocirca totum negotium huc est reductum, ut ex ista aequatione valores
literae $\omega$ elicitur, ubi quidem statim apparet, valorem $\omega = 0$
satisfacere; quia autem hinc oritur $f = \infty$, hincque etiam $k$
infinitum, hoc casu virga nullum plane motum concipiet, sed in quiete
perseverabit. Pro reliquis autem valoribus, quotcunque fuerint inventi,
habebimus
\note{« (119) » en haut, au milieu ; le nombre 79 est à l'encre en haut à
droite. Les signes des équations sont transcrits tels qu'ils ont été lus,
en particulier « $p - q = -2\cos\omega$ » et les deux « $(p - q)$ » de
l'équation qui précède, sans être confrontés au calcul. Le feuillet
s'arrête sur « habebimus ».}


\end{document}
